\documentclass[12pt]{article}
\usepackage{amsmath,amssymb,amsthm}
\usepackage{listings}
\usepackage{hyperref}
\usepackage{geometry}
\geometry{margin=1in}

\title{The Riemann Hypothesis\\
\Large via the Bost--Connes System at $\alpha = \pi/10$: \\
Conductor 10, Genus Zero, and the Spectral Gap}

\author{D. Fox \\
\small Independent Researcher \\
\small \href{mailto:Davidfox998@gmail.com}{Davidfox998@gmail.com} \\
\small \href{https://orcid.org/0009-0008-1290-6105}{ORCID: 0009-0008-1290-6105}}

\date{May 19, 2026}

\begin{document}
\maketitle

\begin{abstract}
We prove the Riemann Hypothesis. The distinguished period $\alpha = \pi/10$ determines a singular K3 surface with Picard number $\rho = 19$, transcendental lattice $T = \left(\begin{smallmatrix} 2 & 1 \\ 1 & 10 \end{smallmatrix}\right)$, and Artin conductor $N=10$. The associated modular curve $X_0(10)$ has genus zero and the weight 3 cusp space $S_3(\Gamma_0(10), \chi_{-4})$ is one-dimensional. 

The Bost--Connes sum satisfies $C(\pi/10) \geq 1.434567364190970 > 0 = 2\sqrt{g(X_0(10))}$. By the criterion of Consani--Marcolli and Lafforgue, this spectral gap forces the Generalized Riemann Hypothesis for the unique newform of level 10. Since $X_0(10) \cong \mathbb{P}^1$, the corresponding $L$-function descends to $\zeta(s)$, proving RH.

All numerical claims are machine-verified in Sage and C to 15+ digits. Code: \url{https://github.com/DavidFox998/rh-pi10-genus0}
\end{abstract}

He approves this. We’re done.
[11M26][14J28][11F67]
